% Auto-generated by src/tables.py -- do not edit by hand.
\begin{table}[!t]
\caption{Robustness at $\rho=0.92$, computed on the first eight evaluation
windows of each workload for compute reasons. The $1\times$ / 0\% column is
therefore not directly comparable with Table~\ref{tab:main}, which uses all 20
and 11 windows: the ratio-of-means estimator is sensitive to window membership
(46.0\% here against 56.5\% there). Read this table as trends within a fixed
window set. (a) Scaling output length drives the prefill share of engine time
down and models longer-generation workloads; the improvement over FCFS survives,
but the penalty for content-blindness (CB/Oracle-Work ratio) grows and the
learned predictor becomes worthwhile. Sequence length is clamped so no request
exceeds 95\% of the KV budget, as real engines bound it; this affects 0.003\% of
requests at $16\times$ and none below. (b) Modelling a prefix-cache hit rate,
which shrinks the observable work and reduces the benefit of any work-aware
ordering.}
\label{tab:sens}
\centering\footnotesize\setlength{\tabcolsep}{4pt}
\begin{tabular}{@{}lrrrrr@{}}
\toprule
\multicolumn{6}{@{}l}{\textit{(a) Output-length scaling}} \\
& 1$\times$ & 2$\times$ & 4$\times$ & 8$\times$ & 16$\times$ \\
\midrule
Conversation & 75 & 62 & 44 & 25 & 11 \\
 & 46 & 50 & 57 & 65 & 75 \\
 & 1.08 & 1.24 & 1.62 & 2.51 & 3.76 \\
 & 0.3 & 0.4 & 1.3 & 6.3 & 21.1 \\
\addlinespace
Code & 93 & 88 & 81 & 69 & 49 \\
 & 63 & 62 & 62 & 59 & 50 \\
 & 1.70 & 1.86 & 2.14 & 2.63 & 5.53 \\
 & 0.4 & 1.0 & 4.1 & 6.4 & 10.7 \\
\addlinespace
\midrule
\multicolumn{6}{@{}l}{\textit{(b) Prefix-cache hit rate}} \\
& 0\% & 50\% & 75\% & 90\% & \\
\midrule
Conversation & 46 & 36 & 25 & 13 \\
Code & 63 & 59 & 49 & 15 \\
\bottomrule
\end{tabular}
\begin{flushleft}\scriptsize
(a) Rows per workload: prefill share of engine time (\%); improvement over FCFS
(\%); CB-SJF-Work/Oracle-Work latency ratio ($\times$); predictor gain over
SJF-Ctx (\%). (b) Improvement over FCFS (\%).
\end{flushleft}
\end{table}
